\documentclass[11pt,a4paper]{article}
\usepackage[top=1.7cm,bottom=1.5cm,left=1.5cm,right=1.5cm,headsep=8pt,footskip=16pt]{geometry}
\usepackage{enumitem}
\usepackage{amssymb}
\usepackage{array}
\usepackage[table]{xcolor}
\usepackage{fancyhdr}
\usepackage{lastpage}
\usepackage{multicol}
\usepackage{polyglossia}
\setdefaultlanguage[calendar=gregorian, numerals=maghrib]{arabic}
\setotherlanguage{english}
\IfFontExistsTF{Amiri}{%
  \newfontfamily\arabicfont[Script=Arabic, Scale=1.08]{Amiri}
}{%
  \newfontfamily\arabicfont[Script=Arabic, Scale=1.08, Path=../fonts/, Extension=.ttf,
    UprightFont=*-Regular, BoldFont=*-Bold, ItalicFont=*-Italic,
    BoldItalicFont=*-BoldItalic]{Amiri}
}
\IfFontExistsTF{AlBattar}{%
  \newfontfamily\khatt[Script=Arabic, Scale=1.15]{AlBattar}
}{%
  \newfontfamily\khatt[Script=Arabic, Scale=1.15, Path=../fonts/,
    BoldFont=ae_AlBattar.ttf]{ae_AlBattar.ttf}
}
\renewcommand{\headrulewidth}{0.8pt}
\renewcommand{\footrulewidth}{0.8pt}
\fancyhf{}
\rhead{\khatt المدرسة العليا للأساتذة ببوزريعة}
\lhead{\khatt ملحقة بومرداس}
\lfoot{\khatt الأستاذ: د. بن عبد الرزاق يوسف}
\cfoot{\thepage/\pageref{LastPage}}
\pagestyle{fancy}
\parindent 0pt
\setlength{\columnsep}{0.6cm}
\setlength{\columnseprule}{0.4pt}
\setlist[enumerate,1]{label=\textbf{\arabic*.}, leftmargin=*, itemsep=5pt, topsep=2pt}
\begin{document}

\linespread{1.08}\selectfont
\begin{center}
{\LARGE\khatt\textbf{الاختبار الخامس والعشرون في الذكاء الاصطناعي}}\\[4pt]
{\normalsize وضعيات إشكالية ــ أسئلة متعددة الاختيارات ــ الفصول من 1 إلى 6}
\end{center}
\vspace{-2pt}\rule{\textwidth}{0.6pt}\\[2pt]
\textbf{التعليمة:} يتضمن الاختبار خمس وضعيات إشكالية. اقرأ كل وضعية بعناية، ثم اختر التحليل أو التصرف الأنسب
من بين الاختيارات الخمسة؛ لكل وضعية إجابة صحيحة واحدة فقط.
اكتب \textbf{حرف} الإجابة الصحيحة في جدول الإجابات في آخر الاختبار.
لكل وضعية 4 نقاط (المجموع 20 نقطة).
\vspace{4pt}

\noindent
{\large\khatt\textbf{الوضعية 1}}\hfill{\small(الفصل 1)}\par\nopagebreak\vspace{3pt}
\setlength{\fboxsep}{7pt}\colorbox{gray!10}{\parbox{\dimexpr\linewidth-14pt}{طلبت أستاذة من أداة توليدية جموع تكسير لعشرين كلمة لنشاط صرفي، فجاءت القائمة منسقة، لكن ثلاثة جموع غير مستعملة في المعاجم، وجمعًا رابعًا صحيحًا لكن لمعنى آخر للكلمة. وترى زميلتها أن الأداة «لا تخطئ في الصرف لأنه قواعد ثابتة». وتريد توزيع القائمة على تلاميذ السنة الرابعة متوسط. ما التقدير الأدق؟}}\par\nopagebreak
\begin{itemize}[leftmargin=1.0cm, labelsep=6pt, itemsep=3pt, topsep=4pt, parsep=0pt]
  \item[\textbf{أ)}] الأداة تولّد صيغًا محتملة بالقياس على أنماط تعلّمتها، فقد تنتج جمعًا قياسيًا غير مسموع، أو تخلط بين معاني الكلمة؛ فتُراجَع الجموع في المعاجم، ويُحدَّد المعنى المقصود في الطلب.
  \item[\textbf{ب)}] الأداة أخطأت لأن الكلمات العشرين كثيرة، ولو طلبت خمس كلمات فقط لجاءت كل الجموع صحيحة بالضرورة، فيكفي تقسيم الطلب دون الرجوع إلى المعاجم أو تحديد المعاني.
  \item[\textbf{ج)}] الأداة لا تعرف الصرف العربي إطلاقًا، وكل ما قدّمته مختلق حتى الجموع الصحيحة منها، فلا يمكن الاستفادة منها في أي نشاط صرفي ولو روجعت مخرجاتها في المعاجم، ثم تبحث عن نشاط آخر.
  \item[\textbf{د)}] تطلب من الأداة التأكد من صحة الجموع، فإن أكدتها اعتمدتها، لأن الأداة تراجع قواعد الصرف في المحاولة الثانية، فتكتشف الجموع غير المسموعة وتحذفها دون حاجة إلى المعاجم.
  \item[\textbf{هـ)}] زميلتها محقة، فالصرف قواعد ثابتة تطبقها الأداة آليًا كالآلة الحاسبة، والجموع غير الموجودة في المعاجم صحيحة قياسًا، والمعاجم هي التي أغفلتها، فتُعتمد القائمة كما هي.
\end{itemize}
\par\vspace{6pt}
\noindent
{\large\khatt\textbf{الوضعية 2}}\hfill{\small(الفصل 3)}\par\nopagebreak\vspace{3pt}
\setlength{\fboxsep}{7pt}\colorbox{gray!10}{\parbox{\dimexpr\linewidth-14pt}{يريد طالب التدرب على الامتحان الشفهي في مقياس الأدب الجاهلي باستعمال أداة ذكاء اصطناعي تؤدي دور الممتحِن. جرّب طلبًا بسيطًا: «اسألني عن الأدب الجاهلي»، فجاءت أسئلة عامة سهلة، وكانت الأداة تمدح كل إجاباته حتى الخاطئة منها. والامتحان بعد أسبوعين، ولا يجد زميلًا يتدرب معه. أي إعداد للطلب أنسب لتدريب فعّال؟}}\par\nopagebreak
\begin{itemize}[leftmargin=1.0cm, labelsep=6pt, itemsep=3pt, topsep=4pt, parsep=0pt]
  \item[\textbf{أ)}] يطلب من الأداة أن تكون «لطيفة ومشجعة دائمًا»، لأن المديح يرفع ثقته بنفسه قبل الامتحان، والتقويم الصريح يُحبطه، فلا حاجة إلى تحديد المحاور أو تصحيح الأخطاء.
  \item[\textbf{ب)}] يحدد الدور والمحاور المقررة ومستوى الصعوبة، ويطلب سؤالًا واحدًا في كل مرة، ثم تقويمًا صريحًا لإجابته بمعايير محددة مع تصحيح الخطأ، ويتحقق من التصحيحات في مراجعه.
  \item[\textbf{ج)}] يحدد الدور والمحاور ويطلب تقويمًا صريحًا، ثم يعتمد كل تصحيحات الأداة دون مراجعتها في المصادر، لأن الأداة في دور الممتحِن لا تخطئ أبدًا في معلومات الأدب الجاهلي.
  \item[\textbf{د)}] يطلب من الأداة كل الأسئلة الممكنة مع إجاباتها النموذجية دفعة واحدة، ثم يحفظ الإجابات، لأن الحفظ أنجع من الحوار في التحضير لأي امتحان شفهي في الأدب.
  \item[\textbf{هـ)}] يطلب من الأداة أن تقلّد أستاذه باسمه وأسلوبه، لأن ذكر اسم الأستاذ يجعل الأداة تعرف أسئلته الحقيقية في الامتحان، فلا حاجة إلى تحديد المحاور أو معايير التقويم.
\end{itemize}
\par\vspace{6pt}
\noindent
{\large\khatt\textbf{الوضعية 3}}\hfill{\small(الفصل 4)}\par\nopagebreak\vspace{3pt}
\setlength{\fboxsep}{7pt}\colorbox{gray!10}{\parbox{\dimexpr\linewidth-14pt}{أعدّت أستاذة بأداة ذكاء اصطناعي مئة جملة لتمارين النحو، فلاحظت أن الجمل تجعل الرجال أطباء ومهندسين وقادة، والنساء في البيت والمطبخ غالبًا، وأن الأسماء كلها تقريبًا من منطقة واحدة. وستوزع التمارين على كل أقسام المتوسطة. وقد اعتادت أن تثق بمخرجات الأداة في التمارين السابقة. ما التصرف الأنسب؟}}\par\nopagebreak
\begin{itemize}[leftmargin=1.0cm, labelsep=6pt, itemsep=3pt, topsep=4pt, parsep=0pt]
  \item[\textbf{أ)}] تحذف كل الجمل التي تتضمن أشخاصًا، وتكتفي بجمل عن الأشياء والحيوانات، لأن الجمل الخالية من البشر لا يمكن أن تنقل أي قيمة ضمنية، فتتخلص من المشكلة نهائيًا.
  \item[\textbf{ب)}] تكتب في الطلب «كن عادلًا» فقط، ثم توزع الجمل الجديدة دون مراجعة، لأن هذه الكلمة تزيل كل صورة نمطية من مخرجات الأداة، فلا حاجة إلى تحديد التنوع المطلوب.
  \item[\textbf{ج)}] توزع الجمل كما هي، لأن هدف التمارين نحوي لا اجتماعي، والتلاميذ يركزون على الإعراب فقط، ولا تؤثر الأدوار المتكررة في الجمل في تصوراتهم عن الرجال والنساء.
  \item[\textbf{د)}] تعكس الأدوار آليًا في كل الجمل، فتجعل كل النساء طبيبات وكل الرجال في المطبخ، لأن العكس التام وحده يصحح الانحياز، ولا حاجة إلى تنويع الأسماء أو السياقات.
  \item[\textbf{هـ)}] تعدّل الطلب بتحديد تنوع الأدوار والأسماء والسياقات، وتراجع الجمل لتصحيح الصور النمطية قبل التوزيع، لأن أمثلة التمارين تنقل قيمًا ضمنية، والانحياز الموروث من البيانات يتسرب إليها.
\end{itemize}
\par\vspace{6pt}
\noindent
{\large\khatt\textbf{الوضعية 4}}\hfill{\small(الفصل 5)}\par\nopagebreak\vspace{3pt}
\setlength{\fboxsep}{7pt}\colorbox{gray!10}{\parbox{\dimexpr\linewidth-14pt}{سجّلت طالبة محاضرة في البلاغة مدتها ساعتان، ثم طلبت من أداة ملخصًا لنصها المفرَّغ لتراجع منه للامتحان. جاء الملخص منظمًا، لكنه حذف الأمثلة والشواهد كلها، ودمج تعريفين متقاربين في تعريف واحد، مع أن الأستاذ ميّز بينهما بإلحاح. والامتحان بعد ثلاثة أيام، والملخص أوضح من نص المحاضرة الطويل. ما الاستعمال الأنسب للملخص؟}}\par\nopagebreak
\begin{itemize}[leftmargin=1.0cm, labelsep=6pt, itemsep=3pt, topsep=4pt, parsep=0pt]
  \item[\textbf{أ)}] تعتمد الملخص وحده في المراجعة، لأن الأمثلة والشواهد تفاصيل لا تُطرح في الامتحان، والتعريف المدمج أوضح من التعريفين المنفصلين، فلا حاجة إلى الرجوع إلى نص المحاضرة.
  \item[\textbf{ب)}] تطلب ملخصًا أقصر من الأول، لأن الملخص القصير يحتفظ تلقائيًا بالتمييزات الدقيقة والشواهد الأساسية، فالحذف والدمج سببهما طول الملخص الأول لا طبيعة التلخيص.
  \item[\textbf{ج)}] تتخلى عن التلخيص وتراجع من التسجيل الصوتي وحده، لأن أي ملخص مكتوب يحذف الأمثلة بالضرورة، ولا يمكن تعديله أو مقابلته بالنص المفرَّغ لاستعادة ما حُذف منه.
  \item[\textbf{د)}] تعدّه خريطة أولية للمحاضرة لا بديلًا عنها، فتعيد إليه التمييز بين التعريفين والشواهد الأساسية من النص المفرَّغ، أو تطلب ملخصًا يحتفظ بالتعريفات حرفيًا، ثم تراجعه بالنص.
  \item[\textbf{هـ)}] تعتمد الملخص بعد أن تضيف إليه الشواهد فقط، وتُبقي التعريف المدمج، لأن الأستاذ يقبل في الامتحان أي تعريف عام، وإلحاحه على التمييز مسألة شكلية لا تؤثر في التقويم.
\end{itemize}
\par\vspace{6pt}
\noindent
{\large\khatt\textbf{الوضعية 5}}\hfill{\small(الفصل 6)}\par\nopagebreak\vspace{3pt}
\setlength{\fboxsep}{7pt}\colorbox{gray!10}{\parbox{\dimexpr\linewidth-14pt}{رفعت أستاذة صورًا لست لوحات خطية إلى أداة متعددة الأنماط لتحديد نوع الخط في كل منها، لدرس عن الخط العربي. حدّدت الأداة أربعًا بدقة، لكنها وصفت لوحة بخط الثلث بأنها بالديواني، وقرأت نص لوحة أخرى قراءة خاطئة جزئيًا. وستعرض اللوحات على تلاميذ يتعرفون إلى الخطوط أول مرة. ما التصرف الأنسب؟}}\par\nopagebreak
\begin{itemize}[leftmargin=1.0cm, labelsep=6pt, itemsep=3pt, topsep=4pt, parsep=0pt]
  \item[\textbf{أ)}] تعتمد تصنيف الأداة كما هو، لأن الثلث والديواني متقاربان، والتلاميذ لا يحتاجون إلى هذا التمييز الدقيق، والقراءة الخاطئة الجزئية لا تؤثر في هدف الدرس عن الخط العربي، فيبقى الدرس بسيطًا ومناسبًا لمستواهم.
  \item[\textbf{ب)}] تعتمد التصنيفات الأربعة الصحيحة دون تحقق، وتحذف اللوحتين الأخريين من الدرس، لأن صحة أربعة تصنيفات تثبت صحتها، ولا حاجة إلى الرجوع إلى أي مرجع في الخط العربي.
  \item[\textbf{ج)}] تراجع تصنيف كل لوحة وقراءتها بالرجوع إلى مرجع في الخط العربي أو مختص، وتصحح الخطأين قبل الدرس، لأن الأداة قد تخلط بين الخطوط المتقاربة، والدقة شرط في درس يعلّم التمييز بينها.
  \item[\textbf{د)}] تتخلى عن الأداة وعن الصور معًا، لأن الأدوات متعددة الأنماط لا تميّز أي نوع من الخطوط، وتصنيفاتها الأربعة الصحيحة جاءت مصادفة، فلا يمكن الاستفادة منها في أي درس.
  \item[\textbf{هـ)}] تطلب من الأداة إعادة التصنيف، وتعتمد التصنيف الذي يتكرر في المحاولتين، لأن التكرار دليل على الصحة، فلا حاجة إلى مرجع أو مختص، حتى في اللوحات التي اختلف فيها التصنيف، كما يفعل كثير من المستخدمين.
\end{itemize}
\par\vspace{6pt}
\begin{center}
{\khatt\textbf{جدول الإجابات}}\\[3pt]
\renewcommand{\arraystretch}{1.6}
\begin{tabular}{|>{\centering\arraybackslash}m{1.9cm}|*{5}{>{\centering\arraybackslash}m{1.2cm}|}}
\hline
\textbf{رقم السؤال} & \textbf{1} & \textbf{2} & \textbf{3} & \textbf{4} & \textbf{5} \\ \hline
\textbf{الإجابة} &  &  &  &  &  \\[8pt] \hline
\end{tabular}

\vspace{8pt}
\renewcommand{\arraystretch}{2.2}
\begin{tabular}{|>{\centering\arraybackslash}m{7.5cm}|>{\centering\arraybackslash}m{3.2cm}|>{\centering\arraybackslash}m{3.2cm}|}
\hline
\textbf{الاسم واللقب} & \textbf{الفوج} & \textbf{العلامة} \\ \hline
 & & \hfill /20 \\ \hline
\end{tabular}
\end{center}
\end{document}
